\chapter{净秽卷}
\chapterkana{穢れ}
\epigraphbox{“澡雪精神。”}{——《庄子·知北游》}

\kaishou{全城写得最薄的一卷，是讲脏的那一卷。\par
一百零四条，管两万年。\par
扫出来的东西堆在城外，堆了两万年，薄薄一层。}

\section{一百零四条}

《章程总目》六卷之中，篇幅最厚的是居止，有二百八十一条；最薄的则是净秽，仅一百零四条。前文只点出这份厚薄次序与常理相左，尚未深究缘由：在一座聚居十余万人的城池里，两万年间管房舍与门禁的条文竟比管清扫的多出两倍有余。若按常理推想，日常琐务越多，规程越该细密；可这里的清扫条文偏偏少得反常。

将这一百零四条拆开来看，其中谈「界」的条文最多，规定何者算界内、何者算出界、越界当如何处置；谈「弃」的次之；谈「封」的最少，仅得九条。三条线索各自指向一个动作：界是划分清楚，弃是移出城外，封是永久封存。合起来便是整套净秽机制的流程，先把某样东西界定为城外之物，送出界外，再加封存。至于那被送出去的东西究竟是什么，一百零四条里没有一条交代。

这套动作倒也并非无迹可寻。静海七号出土的残简里保存着净秽制度的细则，补足了谁来复核、经过哪一道门等细节。\footnote{净秽制度的条文与释文，见《月之都——〈广寒宫残简〉辑佚与考释》论净秽制度的一章。该书处理的是残简，本章用的是章程总目的条目，两者可以互校，条目总数以总目为准。}两边对读，总目如同骨架，残简则是血肉：总目只简略写着「越界者出」，残简则记明越界之人须先带至一道专设的门前，门后有人核对名籍，方才送出城去。细则中最紧要的一条关乎时限：出界者当日不得回环，须在外舍住满三个月。\footnote{同上。三个月这一数字在总目里查不到，只见于残简，两处并不抵牾，可以并用。}这一期限与入城手续中的「先居外舍三月，乃入环」恰好相等。进来要等三个月，出去也要等三个月，中间隔着同一道门。

\section{出界的处置}

第一百零四条是净秽卷的末条，也是全卷最短的一条，仅有两行刻字。

\begin{archive}[内部参考]{档案 № ФВ-1965-09}
净秽卷末条：「凡出界者，不记名，不记数，记其所出之环与时刻。」
枚末另有一行，刀迹不同：「出界之数，岁计卷不载。」
\end{archive}

正文那一句讲的是登记规矩：出界之人既不记姓名，也不记人数，只记从哪一环、在何时出去。枚末那一行刀迹不同的补注则封死了另一条路：连负责统筹物资人口的岁计卷里，也一概不记出界之数。换言之，这套制度在设计之初便刻意回避数量统计，清理了什么、清了多少、清往何处，全部不入总账。之所以如此规定，材料未曾作答，但它留下的后果十分明确：任何想在事后核对总账的人，都无法从账册本身找到答案，只能转去城外寻找旁证。

\section{城外}

账册既空，旁证便只能落在排风口带出城外的堆积物上。由于月面本无大气，城内的空气全靠人工系统维持，城中常年吹拂的定向风便是这套系统的排风。气流经过过滤后，由环壁外侧的开口排出城外，凡随风带出的微粒都会沉降在开口附近。

一九七二年在旧址外侧做过一次浅层取样，掘进两米深，分层记录如下。

\begin{table}[htbp]
\centering
\caption{城外开口附近的堆积}
\begin{tabular}{@{}lll@{}}
\toprule
层位 & 厚度 & 成分\\
\midrule
上层 & 约 30 厘米 & 细尘、凝结物，松散 \\
中层 & 约 1.2 米 & 同上，压实 \\
底层 & 未探到底 & 未见人工物 \\
\bottomrule
\end{tabular}
\tabsource{据一九七二年取样记录。取样点三处，相距十余米，分层一致。}
\end{table}

三处取样点相距十余米，地层剖面却完全一致，说明这是整座城长期排风的结果，并非某一处开口的偶然堆积。排风口的口径与分布位置，也与第一章及《广寒宫残简》所考的环外开口数目相符。\footnote{环外开口共九处，均在环壁下半，与排风走向一致；详见《月之都——〈广寒宫残简〉辑佚与考释》论城形制的一章。}

真正令人费解的是堆积物的成分。一座聚居十一万人的城池运转了两万年，倘若日常清理的是寻常可见的脏污垃圾，这两米厚的堆积里理应夹杂着织物纤维、骨屑、炭粒乃至器物残片，且比例会随着城内起居而波动。可取样记录干干净净，既无纤维也无碎屑，只有单调的细尘与凝结物，也就是排风系统自身运转带出的尘埃。

\section{变化}

堆积的厚度同样在说话。两万年累积一米五，平均每年不到零点一毫米。如此微弱且恒定的速率，似乎与城池规模、人口增减乃至城内发生过何种变故毫无干系。

更直接的线索藏在净秽卷第十二条。这一条并非具体罚则，而是一句总纲式的分类，将需要清除的对象归为三类，每类只用一个字：「秽三者：生、死、易。皆出之。」

将「生」「死」「易」并列一处，它们的共同点呼之欲出。生是生长，死是衰败，易是流转变化；三者无一指代具体的物件，全都是过程。倘若一份清扫规程针对的是有形污物，分类理应按材质与去处来定，诸如布帛、骨殖、灰烬、泥沙各归其类；唯独这一卷按生、死、易立目，分明是在界定事情的性质，而非物件的种类。

面对「城外无物」与「秽分生死易」这两条线索，可以尝试两种截然不同的读法。一种读法认为，清理的对象依然是物质，只是它在城外彻底分解殆尽，故而两万年不留丝毫残渣；另一种读法认为，需要清除的对象原本便非物质，城外取样一无所获反倒成了常态，净秽卷之所以不记数目，也正因为无物可数。

材料本身的倾斜十分明显：净秽卷一百零四条，全篇不谈材质、重量、体积或处理工艺，通篇只围绕界、弃、封三个动作展开。一套处置实体垃圾的规程绝无可能写成这副模样。顺着后一种读法推下去，「秽」所指的其实就是变化本身，包括生长、衰败、腐朽与生死，凡与变化相关的过程皆在此列。净秽制度要清出去的，并非有形的脏污，而是把「变化」二字从这座城里彻底驱逐出去。

一旦接受这一读法，前面几章留下的几处疑团便迎刃而解。章程里没有医卷也没有育卷，只因城内无人会病、亦无人生育；净秽卷薄得只剩一百零四条，只因要清除的对象本无形质，无须按物细分；至于那一千二百四十七条章程两万年不改一字，也正是因为判断轻重缓急的那套尺度随着变化一并停摆了。

\section{限度}

平心而论，这一章撑起全部推论的柱子只有两根：一百零四条简略的条文，以及三处两米深的浅层取样。两样材料都算不上厚重，条文刻意不记数，取样又只覆盖了三个点；倘若读者向笔者索要一件「被清出来的东西」作实物证明，本章一件也拿不出来，因为按这一读法，那东西本来就不该有实物留存。

笔者能担保的，仅限于当前材料确实指向「非物」这一侧，且这一猜想恰好能与全书其余各章严丝合缝地对上。当然，各处对得上并不自动等于事实果真如此，它只说明目前尚未出现直接推翻它的反证。若这一前提暂且成立，城内不老、不病、不育、技术停滞等种种反常便无需再逐一找补，它们全都是同一件事的不同侧面。至于被清出去的「变化」究竟去往何处、又由谁经手处理，则是下一章必须面对的难题。
